% Vietnamese translation for OI Public Library
% Source-File: IOI中国国家候选队论文/2006/陈首元/陈首元.ppt
\documentclass[11pt]{article}
\usepackage{oipl-vn}

\newcommand{\Oh}{\mathcal{O}}

\title{Cây động\\{\large Ghi chú theo slide}}
\author{Chen Shouyuan}
\date{2006}

\begin{document}
\maketitle

\origtitle{Dynamic trees}
\sourcepath{IOI China candidate team papers / 2006 / Chen Shouyuan / Chen Shouyuan.ppt}

\section{Phạm vi}

Bộ slide đi kèm bài viết về cấu trúc cây động để duy trì một rừng có trọng
số. Ghi chú này dựa trên bản bài viết đã dịch.

\section{Các thao tác cần hỗ trợ}

Cây động duy trì một rừng và hỗ trợ nối hai cây, cắt một cạnh, đổi gốc, tìm
gốc của đỉnh, cập nhật trọng số trên đường từ gốc tới đỉnh, và truy vấn cạnh
có trọng số nhỏ nhất trên đường đó. Cách lưu cha trực tiếp làm một số thao tác
nhanh, nhưng các truy vấn theo đường có thể tốn \(\Oh(n)\).

\section{Cạnh thực và cạnh ảo}

Ý tưởng chính là chia cạnh thành cạnh thực và cạnh ảo. Mỗi đỉnh có nhiều nhất
một cạnh thực đi xuống, nên các cạnh thực tạo thành các đường rời nhau. Thay
vì thao tác trên toàn bộ cây, ta thao tác trên các đường này và thay đổi cách
phân rã khi cần.

\section{\texttt{expose}}

Thao tác trung tâm là đưa đường từ gốc tới một đỉnh thành một đường thực.
Trong quá trình đó, một số cạnh ảo được đổi thành cạnh thực và ngược lại. Nếu
cấu trúc lưu mỗi đường hỗ trợ tách, ghép, đảo và cập nhật lười, các thao tác
cây động đạt \(\Oh(\log n)\) khấu hao.

\section{Ứng dụng}

Cây động hữu ích trong các bài cần duy trì thông tin trên rừng thay đổi theo
thời gian, đặc biệt là thuật toán tối ưu hóa mạng và các biến thể của cây
khung động. Bài viết nhấn mạnh rằng sức mạnh của cấu trúc đến từ việc biến
truy vấn cây thành truy vấn trên một số ít đường.

\end{document}
